\documentclass{infthesis}
\usepackage[utf8]{inputenc}
\usepackage[T1]{fontenc}
\usepackage{geometry}
\usepackage{graphicx}
\usepackage{float}
\usepackage{hyperref}
\usepackage[english,indonesian]{babel}

\title{Optimasi Pemrosesan Template LaTeX Menggunakan Pendekatan Berbasis AI}
\author{Academic Support Agent}
\date{\today}

\begin{document}
\maketitle

\begin{abstract}

\end{abstract}

\chapter{PENDAHULUAN}


\chapter{TINJAUAN PUSTAKA}


% Original content (commented out):
% %%%%%%%%%%%%%%%%%%%%%%%%
% Sample use of the infthesis class to prepare an MSc thesis.
% This can be used as a template to produce your own thesis.
% Date: June 2019
%
%
% The first line specifies style options for taught MSc.
% You should add a final option specifying your degree.
% *Do not* change or add any other options.
%
% So, pick one of the following:
% \documentclass[msc,deptreport,adi]{infthesis}     % Adv Design Inf
% \documentclass[msc,deptreport,ai]{infthesis}      % AI
% \documentclass[msc,deptreport,cogsci]{infthesis}  % Cognitive Sci
% \documentclass[msc,deptreport,cs]{infthesis}      % Computer Sci
% \documentclass[msc,deptreport,cyber]{infthesis}   % Cyber Sec
% \documentclass[msc,deptreport,datasci]{infthesis} % Data Sci
% \documentclass[msc,deptreport,di]{infthesis}      % Design Inf
% \documentclass[msc,deptreport,inf]{infthesis}     % Informatics
%%%%%%%%%%%%%%%%%%%%%%%%

\documentclass[msc,deptreport,slp]{infthesis} % Do not change except to add your degree (see above).

\begin{document}
\begin{preliminary}

\title{This is the Project Title}

\author{Your Exam Number\\word count}

\abstract{
This is the infthesis template modified to give the right affliations for the MSc SLP. 

  This skeleton demonstrates how to use the \texttt{infthesis} style. 
  Note \textbf{there isn't a page limit for SLP dissertations but there is a word limit (8000 words)}. 
  
  The file \texttt{skeleton.tex} generates this document and
  can be used as a starting point for your thesis. The abstract should
  summarise your report and fit in the space on the first page.
}

\maketitle

\section*{Acknowledgements}
Any acknowledgements go here. 

\tableofcontents
\end{preliminary}


\chapter{Introduction}

The preliminary material of your report should contain:
\begin{itemize}
\item
The title page.
\item
An abstract page.
\item
Optionally an acknowledgements page.
\item
The table of contents.
\end{itemize}

As in this example \texttt{skeleton.tex}, the above material should be
included between:
\begin{verbatim}
\begin{preliminary}
    ...
\end{preliminary}
\end{verbatim}
This style file uses roman numeral page numbers for the preliminary material.

The main content of the dissertation, starting with the first chapter,
starts with page~1. \emph{\textbf{SLP dissertations have no page limit, but there is an 8000 word limit.}}

The report then contains a bibliography and any appendices, which may go beyond
the word limit. The appendices are only for any supporting material that's important to
go on record. However, \textbf{you cannot assume markers of dissertations will read them.}



\section{Using Sections}

Divide your chapters into sub-parts as appropriate.

\section{Citations}

Citations (such as \cite{P1} or \cite{P2}) can be generated using
\texttt{BibTeX}. For more advanced usage, the \texttt{natbib} package is
recommended. You could also consider the newer \texttt{biblatex} system.

These examples use a numerical style  (e.g. [1]) or 
(Author, Date) format for inline citations.\footnote{Catherine has a strong preference for (Author, Date) as a marker.  Also, try to avoid using footnotes too much. If you want the marker to read it put it in the main text!}

\chapter{Background}

A dissertation usually contains several chapters.  One should cover the necessary background material to understand why you are doing this study and motivate the methods you choose to use. 

\chapter{Method/Experimental Setup}

You may want to use a different structure depending on the focus of your dissertation. For example, it may make sense to have a chapter devoted specifically to data or architectures.  

It's especially important in ``methods'' chapters to think about reproducibility. If you were to give this dissertation to another SLP student in a year's time, would they be able to reproduce your results? 

\chapter{Results}

You should definitely have a results chapter.  You can also include specific discussion of your results here.

It's a good to link your results to your specific hypotheses/research questions as your present/discuss them. 

\chapter{General Discussion}

A general discussion chapter will often strengthen your dissertation.  It's a good place to discussion your results overall, and how they tie into the issues raised in your introduction and background sections.  It's also good to identify potential areas that could be improved in terms of study design and general future work (nobody expects you to solve everything in 2.5 months!).  Try to identify specific improvements (e.g. what you could do if you had a few more months) and longer term implications.

\chapter{Conclusions}

A good conclusion will concisely summarise/highlight what you found, how it relates to your research questions/hypotheses, and what implications it has for other research.


\bibliographystyle{apalike}
\bibliography{mybibfile}

%% You can include appendices like this:
% \appendix
% 
% \chapter{First appendix}
% 
% \section{First section}
% 
% Markers do not have to consider appendices. Make sure that your contributions
% are made clear in the main body of the dissertation (within the page limit).

\appendix

\chapter{Additional Materials}

You can add materials in an appendix but don't assume that the marker will read it.  It's a good place to include listening test materials, or tables of dev set results, e.g. for hyper parameter tuning.  Stuff that the marker doesn't really need to know for your overall argument but might like to check for reproducibility. 

\end{document}


\end{document}